% Carta de presentación para el envío del artículo 2 a Public Transport (Springer). Compilar con: latexmk -pdf carta-editor.tex
\documentclass[11pt,a4paper]{letter}
\usepackage[T1]{fontenc}
\usepackage[utf8]{inputenc}
\usepackage{lmodern}
\usepackage{microtype}
\usepackage[a4paper,margin=25mm]{geometry}
\usepackage[hidelinks]{hyperref}
\setlength{\parskip}{0.6em}

\signature{Yeray Garrido\\Independent researcher, Bilbao (Biscay), Spain\\ORCID
\href{https://orcid.org/0009-0008-0459-0294}{0009-0008-0459-0294}}
\address{Yeray Garrido\\Bilbao, Biscay, Spain}
\date{\today}

\begin{document}
\begin{letter}{The Editors-in-Chief\\\emph{Public Transport -- Planning and Operations}\\Springer}

\opening{Dear Editors,}

I am pleased to submit the manuscript \emph{``Accuracy of open-data versus operator real-time bus arrival predictions:
a paired evaluation of the Bide+ method in Biscay, Spain''} for consideration as an original research article in
\emph{Public Transport}.

Real-time passenger information is now standard, but the accuracy of the predictions that operators actually show
passengers is rarely measured, and almost never against an alternative on the same vehicles and at the same moments.
The manuscript does this for Bizkaibus, the interurban bus network of Biscay. For one week, the official waiting times
of 60 stops were recorded every minute and compared, bus by bus, with the published timetable and with the arrivals
predicted by the Bide+ method from the operator's open data alone (vehicle positions and timetable). Each announcement
was matched to its vehicle and confirmed with the distance reported by the official service, and the actual passage was
reconstructed in two independent ways.

Across 555,744 paired predictions on six days, the timetable was within two minutes of the actual passage in 16\% of
cases, the official real-time information in 39\% and the open-data method in 64\%. The two real-time sources were
equally accurate below three minutes, and the open-data method was more often within two minutes from three minutes on,
at 57 of 59 stops. The manuscript reports the limits of the comparison openly: coverage, the shared ground truth,
matching uncertainty and the larger mean error of the open-data method below six minutes.

I believe the manuscript fits the scope of \emph{Public Transport} because it concerns the information management and
operation of a real public transport network, and because its main implication is practical: operators that publish
vehicle positions, and also their own predictions, allow the quality of their passenger information to be audited
independently and continuously. The paired protocol, the capture program, the data and the analysis are openly
available, so that the comparison can be repeated on other networks.

The manuscript has been posted as a preprint on Zenodo
(\href{https://doi.org/10.5281/zenodo.24286027}{doi:10.5281/zenodo.24286027}). The method it evaluates is described in
a separate preprint (\href{https://doi.org/10.5281/zenodo.24285540}{doi:10.5281/zenodo.24285540}) and in a technical
report (\href{https://doi.org/10.5281/zenodo.24041572}{doi:10.5281/zenodo.24041572}). Apart from these preprints, the
work has not been published and is not under consideration elsewhere.

As stated in the manuscript, I designed the Bide+ method and develop the free passenger information app that uses it,
which accepts voluntary donations to pay for its server; I receive no other income from it and have no relationship with
the operator. To limit the effect of this interest, the method was frozen before the test week and all data and code are
public.

Thank you for considering this manuscript. I look forward to hearing from you.

\closing{Yours sincerely,}

\end{letter}
\end{document}
